% ===== 1. Introduction (DRAFTABLE) ==========================================
\section{Introduction}
\label{sec:intro}

Differentiable rasterization optimizes an explicit 3D representation directly
against multi-view images by backpropagating a photometric loss through the
rasterizer~\cite{loper2014opendr,kato2018neural,liu2019softras,laine2020nvdiffrast}.
A \emph{triangle soup}---an unstructured collection of triangles
with no shared connectivity---is an attractive such representation: it is cheap
to render, places no manifold constraint on the geometry, and supports
\emph{adaptive resampling}, a periodic in-loop step that prunes
poorly-contributing triangles and respawns new ones to redistribute resolution
where the image residual demands it, in the spirit of the adaptive
densify-and-prune loops of explicit-primitive and mesh-evolution
pipelines~\cite{kerbl2023gaussians,held2025trianglesplatting,nicolet2021largesteps,palfinger2022remeshing}.
% NOTE: the DiffSoup renderer/trainer this instantiates is our own (unpublished)
% system; add its citation here once it is public.

Such pipelines are tuned to minimize geometric error, and geometric error is
exactly what they reduce. But geometry and \emph{topology} are different: a
reconstruction can be everywhere close to the target surface yet wrong in genus,
in the number of connected components, or in whether an enclosed void survives.
Our first contribution is to show that the standard pipeline is
\emph{topologically blind} in a precise, measurable sense. We construct matched
pairs of candidate reconstructions whose Chamfer distance to the ground truth is
\emph{equal by construction} but whose topology differs, and show that
Chamfer---and even Hausdorff---cannot separate them, whereas the
bottleneck distance between persistence diagrams can
(Section~\ref{sec:results-blindness}). We further observe that a topological
bias injected only at \emph{initialization} is erased: the first resampling step
washes it out, so the final topology is independent of the initial bias.
% src(qualitative, washout/blindness): D:\Project\CG-Soup-Topology\figures\summary.md
% TODO(human): verify the one-line phrasing of the washout claim matches how you
% want to state Phase-1 finding F1 (init-bias erased by first resample ~step 100).

Two consequences follow, and they frame the paper. First, correcting topology
requires guidance applied \emph{in the loop}, not at initialization. Second,
detecting whether topology is correct requires a metric that sees it---a
persistence-based measure rather than a point-to-point distance.

We introduce \emph{topology-aware adaptive resampling}
(Section~\ref{sec:method}): a precomputed importance field that localizes the
target's persistent features and biases (i) where new triangles respawn and
(ii) which triangles are protected from pruning. The method is budget-neutral
(the triangle count is held fixed) and reduces \emph{exactly} to the baseline
when its single strength knob is set to zero. Within this one mechanism we vary a
\emph{concentrate$\leftrightarrow$spread} scale on the importance field, which
lets us ask not merely whether topological guidance helps, but \emph{what spatial
form} the guidance should take.

Our central finding concerns the \emph{spatial form} of that guidance. A
topological prior helps voids (H2), but chiefly through \emph{width}: a spread prior
lowers bottleneck-to-target by 33--53\% across spheres, cubes, and capped cylinders,
yet a width-matched non-topological control recovers most of it---so the void gain is
mainly a spread effect, with only a small, shape-dependent topological residual
(Section~\ref{sec:results-h2}).
% src(33--53% = B4 vs B0: sphere 33 / cube 53 / cylinder 34): D:\Project\CG-Soup-for-Digital-Dentistry\output\synth\h2_unified\report\results.json
But the natural hypothesis---that high-dimensional features want a
\emph{concentrating} prior and low-dimensional ones a \emph{spreading} one, a
dimensional crossover---turns out to be \emph{false}. Comparing the two priors
directly (Section~\ref{sec:results-crossover}), a spread prior is at least as good
as concentrating at every dimension and strictly better for voids and loops;
concentrating on a one-dimensional loop actively backfires, over-tessellating it
into phantom handles. Across dimensions the benefit tracks \emph{width}, not
topological placement---a width-matched non-topological prior does at least as well
(better, on the loop)---so the operative takeaway is to \emph{spread} resampling
resolution, not to target topology per se. The unifying reason (Section~\ref{sec:mechanism}) is that a
feature's death-simplex is only a localized witness of its support, so spreading
importance over that support is the right bias---while the \emph{harm} from
concentrating is what depends on dimension. This turns ``does topology help?''
into a prescription for how to shape the prior.
% src(spread>concentrate): D:\Project\CG-Soup-for-Digital-Dentistry\output\synth\crossover\crossover_report\results.json
